% The next line makes it so it compiles the main file if I hit the compile button
% when editing this file in TeXworks.
% !TEX root = ../AIDMA.tex
\section{Problems}

\begin{pro}
Draw a Venn diagram showing $A\cap(B\cup C)$, where $A$, $B$, and $C$
are sets.
\end{pro}


%Now an eval example.
%\begin{pro}
%Use a set containment proof to prove that if $A$ and $B$ are sets, then $A-B = A\cup\overline{B}$.
%\end{pro}

\begin{pro}
Assume $A$, $B$, and $C$ are sets.
Prove each of the following set identities using a set containment proof based on the basic definitions
of $\cap$, $\cup$, etc. (see examples~\ref{setpf1},~\ref{setpf2}, and~\ref{setpf3}).
\begin{lenum}
\item
$(A\cap B\cap C) \subseteq (A\cap B)$.
\item 
$A\cap B \subseteq A\cup B$.
\item 
$(A\cup B)\setminus(A\cap B) = (A\setminus B)\cup(B\setminus A)$.
\item 
$(A-B)\setminus C \subseteq A\setminus C$.
%\item 
%$A\cup(B\cap C)=(A\cup B)\cap (A\cup C)$.
\end{lenum}
\end{pro}

\begin{pro}
Prove each of the following set identities using a set containment proof based on the basic definitions
of $\cap$, $\cup$, etc. (see examples~\ref{setpf1},~\ref{setpf2}, and~\ref{setpf3}).
%I don't have any examples of the following in the book.  Oops.
%{\em and} a proof that uses other set identities.
\begin{lenum}
\item
$A\cup(A\cap B)=A$.
\item
$A\cup (B\cap C) = (A\cup B)\cap(A\cup C)$
\item
$(A\setminus B)\setminus C=(A\setminus C)\setminus (B\setminus C)$.
\item
$\overline{A\cup(B\cap C)}=(\overline{C}\cup\overline{B})\cap\overline{A}$. (This one is a little tricky.)
\end{lenum}
\end{pro}


\begin{pro} Rusty has $20$ marbles of different colors: black, blue, green,  and yellow. 
Seventeen of the marbles are not
green, five are black, and $12$ are not yellow. How many blue
marbles does he have?
\end{pro}

\begin{pro}
You need to settle an argument between your boss (who can fire you) and your professor (who can fail you).  They are trying to decide who to invite to the Young Accountants Volleyball League.  They want to invite freshmen who are studying accounting and are at least 6 feet tall.  They have a list of all students.
\begin{lenum}
\item
Your boss says they should make a list of all freshmen, a list of all accounting majors, and a list of everyone at least 6 feet tall.  They should then combine the lists (removing duplicates) and invite those on the combined list. 
Is he correct?  Explain.  If he is not correct, describe in the simplest possible terms who ends up on his guest list.
\item
Your professor says they should make a list of everyone who is not a freshman, a list of everyone who does not do accounting, and a list of everyone who is under 6 feet tall. They should make a fourth list that contains everyone who is on all three of the prior lists.  Finally, they should remove from the original list everyone on this fourth list, and invite the remaining students.
Is he correct?  Explain.  If he is not correct, describe in the simplest possible terms who ends up on his guest list.
\item
Give a simple description of how the guest list should be created.
\end{lenum}
\end{pro}

\begin{pro}
Using words, explain what $53 \bmod 7 = 4$ means.
\end{pro}

\begin{pro}
Let $a$ be a integer such $a\leq 17/3$. What can you say about $a$?  Prove your claim.
\end{pro}

\newpage
\begin{pro}
Compute each of the following. If 2 answers are possible, give both.
\begin{multicols}{4}
\begin{lenum}
\item $23 \bmod 10$
\item $-14 \bmod 10$
\item $8675309 \bmod 10$
\item $3 \bmod 5$
\item $34 \bmod 5$
\item $-8 \bmod 5$
\item $13 \bmod 12$
\item $144 \bmod 12$
\item $7 \bmod 12$
\item $-7 \bmod 12$
\item $3 \bmod 2$
\item $-3 \bmod 2$
\end{lenum}
\end{multicols}
\end{pro}

\begin{pro}
Compute each of the following.
\begin{multicols}{4}
\begin{lenum}
\item $\lfloor 2.72\rfloor$
\item $\lceil 2.72\rceil$
\item $\lfloor 3.1415\rfloor$
\item $\lceil 3.1415\rceil$
\item $\lfloor\lceil 3.1415\rceil\rfloor$
\item $\lceil\lfloor 3.1415\rfloor\rceil$
\item $\lfloor 3/2\rfloor$
\item $\lfloor 5/4\rfloor$
\item $\lfloor 8.675309\rfloor \bmod 3$
\item $\lceil 8.675309\rceil \bmod 3$
\end{lenum}
\end{multicols}
\end{pro}

\begin{pro}
Prove or disprove: If $a$ is a real number and $n$ is an integer, then $\lfloor a \rfloor \bmod n = \lfloor a\bmod n \rfloor$.
\end{pro}

\begin{pro}
Recall that 
$a\equiv b \pmod n$ iff $a-b=kn$ for some integer $k$.
Use this definition of congruence modulo $n$ 
to prove Theorem \ref{thm:congruence}.  
(Note: This is an if and only if proof, so you need to prove both ways.)
\end{pro}

\begin{pro}
Let $a,b\in\reals$, $a\neq 0$, and define $f:\reals\to\reals$ by $f(x)=ax+b$.
Prove that $f$ is one-to-one and onto.
\end{pro}

\begin{pro}
Let $a$ and $b$ be real numbers with $a\neq 0$.
Define $f:\reals\to\reals$ by $f(x)=a\, x + b$.
Show that $f$ is invertible. Then find $f^{-1}$.
\end{pro}

\begin{pro}
Prove or disprove: if $a$, $b$, and $c$ are real numbers with $a\neq 0$,
then the function $f(x)=a\,x^2+b\, x + c$ is invertible.
\end{pro}

\begin{pro} 
Prove that if $f$ and $g$ are onto, then $f\circ g$ is also onto.
\end{pro}

\begin{pro}
Let $f(x)=x+\lfloor x\rfloor$ be a function on $\reals$. (This one is a little tricky.)
\begin{enumerate}[(a)]
\item Prove or disprove that $f$ is one-to-one.
\item Prove or disprove that $f$ is onto.
\item Prove or disprove that $f$ is invertible.
\end{enumerate}
\end{pro}

\begin{pro}
Find the inverse of the function $f(x)=x^3+1$ over the real numbers.
\end{pro}

\begin{pro}
Let $f$ be the function on $\BBZ^{+}$ that maps  $x$ to the number of
bits required to represent $x$ in binary.  For instance, $f(1)=1$, $f(2)=2$, 
$f(3)=2$, $f(4)=3$, $f(10)=4$, etc.  
Hint:  The number $2^n$ requires $n+1$ bits to represent (a single $1$ followed
by $n$ zeros).  You may be able to use this fact in one of your proofs.
\begin{enumerate}[(a)]
\item Prove or disprove that $f$ is one-to-one.
\item Prove or disprove that $f$ is onto.
\item Prove or disprove that $f$ is invertible.
\end{enumerate}
\end{pro}

\newpage
\begin{pro}
\item 
Consider the relation 
$R=\{(1,2), (1,3), (3,5), (2,2), (5,5), (5,3), (2,1), (3,1)  \}$ on set $\{1,2,3,4,5\}$.
Is $R$ reflexive? symmetric? anti-symmetric? transitive?
an equivalence relation? a partial order?
\end{pro}

\begin{pro}
Let $X$ be the set of all people.
Which of the following are equivalence relations?  Prove it.
\begin{lenum}
\item
$R_1=\{(a,b)\in X^2 | a \mbox{ and } b \mbox{ are the same height} \}$
\item
$R_2=\{(a,b)\in X^2 | a \mbox{ is taller than } b \}$
\item
$R_3=\{(a,b)\in X^2 | a \mbox{ is at least as tall as } b \}$
\item
$R_4=\{(a,b)\in X^2 | a \mbox{ and } b \mbox{ have the same last name} \}$
\item
$R_5=\{(a,b)\in X^2 | a \mbox{ has the same kind of pet as } b \}$
\end{lenum}
\end{pro}

\begin{pro}
Repeat the previous problem, but which are partial orders?  Prove it.
\end{pro}

\begin{pro}
Define three different equivalence relations on the set of all TV shows.
For each, give examples of the equivalence classes, 
including one representative from each.  Prove that each is an 
equivalence relation.
\end{pro}

\begin{pro}
Define a relation on the set of all Movies that is {\em not} an equivalence relation.
\end{pro}

\begin{pro}
Let $A=\{1, 2, \ldots, n\}$.  Let $R$ be the relation on $P(A)$ (the power set of $A$)
such that $a,b\in P(A)$ are related iff $|a|=|b|$.  Prove that $R$ is an equivalence relation.
What are the equivalence classes of $R$?
\end{pro}


